% ==============================================================================
% CEBE Interdisciplinary Fellowship Programme -- Call for PhD and PostDoc projects
% Project proposal template
%
% Formatting requirement from the call: Times New Roman, 12 pt,
% single-spaced text, 2.5 cm margins.
%
% Font choice:
%   This template uses `mathptmx`, a Times-metric-compatible font that
%   ships with every TeX distribution (no extra font files needed, works
%   the same on Windows/MiKTeX, Overleaf, and Linux TeX Live). It is
%   visually indistinguishable from Times New Roman and satisfies "Times
%   New Roman, 12pt" in practice.
%
%   If you would rather use the literal Times New Roman TTF installed on
%   this Windows machine, compile with XeLaTeX or LuaLaTeX instead and
%   swap the font block below for:
%       \usepackage{fontspec}
%       \setmainfont{Times New Roman}
%   (comment out \usepackage{mathptmx} in that case).
%
% Build: pdflatex -> bibtex -> pdflatex -> pdflatex   (or just `latexmk -pdf main.tex`)
% ==============================================================================

\documentclass[12pt,a4paper]{article}

% ---- Page geometry: 2.5 cm margins on all sides ----
\usepackage[a4paper,margin=2.5cm]{geometry}

% ---- Font: Times-compatible, 12pt (set via documentclass option above) ----
\usepackage[T1]{fontenc}
\usepackage[utf8]{inputenc}
\usepackage{mathptmx}

% ---- Single-spaced body text (this is already LaTeX's default, set explicitly) ----
\usepackage{setspace}
\singlespacing

% ---- Useful packages ----
\usepackage{graphicx}      % \includegraphics for a CEBE/university logo on the frontpage
\usepackage{enumitem}      % tighter, controllable lists
\usepackage{booktabs}      % clean tables (applicant list)
\usepackage[numbers,sort&compress]{natbib}  % references; bibliography style set below
\usepackage{xcolor}
\usepackage[colorlinks=true,linkcolor=blue,citecolor=blue,urlcolor=blue]{hyperref} % load near-last
\usepackage {ragged2e}


% ---- Section numbering to mirror the call's own numbering (1, 2, 2.1 ... , 3, 4) ----
\setcounter{secnumdepth}{2}
\setcounter{tocdepth}{2}

% ---- Lightweight page-limit tracker -----------------------------------------
% Prints "PAGEMARK: <label> -- page N" to the compile log at each call. Grep
% the log (main.log) for "PAGEMARK" after compiling to check where each
% page-limited section starts/ends against the call's stated maxima:
%   Project Summary               max 1/2 page
%   Project Description           max 4 pages (excluding references)
%   Project relevance to CEBE     max 1/2 page
%   Sustainability goals          max 1/2 page
\newcommand{\pagemark}[1]{\typeout{PAGEMARK: #1 -- page \thepage}}

\newcommand{\projectname}{\textbf{HAI-BRI }} % iDT, HAI-TWIN
\newcommand{\asset}{bridge case study }

% ---- Expected-length annotations -------------------------------------------
% Prints the call's stated page limit right under each section heading, so
% it stays visible while drafting. Set \templatenotesfalse (just below) to
% hide these before the final submission build -- they are not part of the
% proposal content.
\newif\iftemplatenotes
\templatenotestrue
\newcommand{\lengthnote}[1]{%
  \iftemplatenotes
    {\small\itshape\color{gray!70!black} [Target length: #1]}\par\vspace{4pt}
  \fi
}

\begin{document}

% ==============================================================================
% FRONTPAGE
% Required content: title of project, name and affiliation of applicants,
% table of contents.
% ==============================================================================
\begin{titlepage}
\centering
\vspace*{1.5cm}

% Uncomment and point at a logo file if you want the CEBE / university mark here:
% \includegraphics[width=4cm]{cebe-logo.png}\par\vspace{1cm}

{\large CEBE Interdisciplinary Fellowship Programme}\par
{\large Call for PhD and PostDoc Projects}\par
\vspace{1.5cm}

{\Huge\bfseries Human-AI Collaborative Workflows for Immersive Digital Twins in Bridge Monitoring and Inspection}\par
\vspace{0.5cm}
{\large PhD project}\par

\vspace{2cm}

\begin{tabular}{@{}p{3.2cm}p{4.5cm}p{6.5cm}@{}}
\toprule
\textbf{Role} & \textbf{Name} & \textbf{Affiliation} \\
\midrule
Principal Investigator & Assoc. Prof. Giuseppe Abbiati & Department of Civil and Architectural Engineering, Aarhus University \\
Co-supervisor & Assist. Prof.  Jhonattan Martinez & Department of Civil and Architectural Engineering, Aarhus University \\
Co-supervisor & Dr. Carolin Reichherzer & Department of Civil, Environmental and Geometic Engineering, ETH Zurich \\
Co-supervisor & Prof. Ueli Angst & Department of Civil, Environmental and Geometic Engineering, ETH Zurich \\
Candidate & Lorenzo Loyola & Department of Civil and Architectural Engineering, Aarhus University \\
\bottomrule
\end{tabular}

\vspace{2cm}
%\date{}
Submission date: 2026-09-30

\end{titlepage}

\newpage
\tableofcontents
\newpage

% ==============================================================================
% 1. PROJECT SUMMARY  (max 1/2 page)
% ==============================================================================
\pagemark{Project Summary -- start}
\section{Project Summary}
\lengthnote{max 1/2 page (per the call)}

\textcolor{blue}{Critical infrastructure is increasingly affected by aging, climate exposure, water ingress, and limited maintenance resources, while bridge inspection generates growing volumes of heterogeneous information from Uncrewed Aerial Systems (UAS), sensors, inspection reports, and 3D models. The emerging challenge is to support engineers what the combined evidence means, how certain that interpretation is, and what steps should be taken next, in collaborative workflows that combine AI with engineering expertise. This PhD will develop and validate a human-centered situational-awareness framework for bridge Structural Health Monitoring (SHM). The framework will integrate UAS imagery, sensor measurements, structural information, inspection histories, environmental observations, and engineering knowledge through a Digital Twin (DT) and a Knowledge-grounded (KG). AI based on Large Language Models (LLMs) and Vision-Language Models (VLMs) will support interpretation of this evidence by identifying relevant relationships, plausible deterioration mechanisms, uncertainty, and information gaps. The scientific contribution is a new human–AI decision paradigm in which multimodal evidence is semantically contextualized, transformed into an evolving assessment of the infrastructure situation, interpreted through AI-supported reasoning, challenged by the engineer, and converted into inspection or maintenance decisions. The project will focus on bridge deterioration associated with water ingress, expansion joints, drainage deficiencies, moisture, cracking, and corrosion-related symptoms. The framework will be validated experimentally by comparing conventional engineering workflows, AI assistance without semantic grounding, and the proposed knowledge-grounded human–AI workflow. Collaboration with Aarhus Municipality will provide access to diverse infrastructure case studies across Aarhus, enabling validation under real operational conditions. An international \asset in Salvador, Bahia, Brazil, supported by the Federal University of Bahia, will complement the Danish cases and investigate the framework’s transferability to contexts with different climatic conditions, data availability, monitoring capabilities, and resource constraints. This comparison will also assess the approach's applicability in developing countries, where infrastructure owners may face limited monitoring technologies, incomplete historical data, and constrained inspection and maintenance resources.}

\pagemark{Project Summary -- end}

% ==============================================================================
% 2. PROJECT DESCRIPTION  (max 4 pages, excluding references)
% Required subsection headings -- keep exactly as below.
% ==============================================================================
\pagemark{Project Description -- start}
\section{Project Description}
\lengthnote{max 4 pages total for this section (all subsections combined), excluding references (per the call)}
\lengthnote{Include a clear description of the visions and goals, the distinguishing features, and foci; Include a motivation for why a doctoral student or postdoc is most 
appropriate for the interdisciplinary project}
\subsection{Motivation, Significance and Scientific Challenges}
%\subsubsection{Motivation}
Deterioration of bridge infrastructures is a critical issue in the west, where such assets, in large part, already exceed their nominal life. Europe has a very dense network of transportation infrastructures. The prioritization of what to repair and what to demolish is a critical challenge. The catastrophic collapse of the Morandi Bridge (Genova, Italy, 2018)indicates that bridges in poor condition may still be operated. The situation in global south is characterized by younger infrastructures which, however, due to underfunded maintenance programs are likely to experience fast deterioration. SHM has been implemented extensively in the last 20 years along with the large diffusion of cost-effective sensing technology, broadband internet access, and connectivity through IoT. More recently, UASs have become affordable and easy to operate, facilitating the capture of high-resolution image data from hard-to-access structural elements. However, engineers must still combine these heterogeneous sources to determine whether an observed condition is significant, whether additional inspection is needed, and which intervention to prioritize~\cite{hoskere_unified_2025}. The overwhelming amount of data, along with the shortage of qualified personnel, renders this a challenging task. \par
%
%The primary contribution of this work is to integrate them within one human-AI collaborative workflow for immersive digital twins for bridge monitoring and inspection in which AI reasoning and human engineering judgment have defined and complementary roles. The workflow will support evidence retrieval, contextualization, alternative interpretations, and identification of information gaps, while humans remain responsible for consequential maintenance decisions.
%Nevertheless, improved data acquisition does not automatically lead to better engineering decisions. Inspection information remains widely distributed across heterogenous data sources (images, point clouds, drawings, reports, sensor measurements, historical observations, and expert knowledge) .
%
%\subsubsection{Significance}
%
The project's vision is to develop a human-centered DT that supports engineers in designing and implementing inspection campaigns. A key aspect is enabling engineers to build situational awareness based on the data collected and build the foundation on how to proceed with inspection or intervention. AI models may provide an opportunity for engineers to rapidly access and interpret a variety of data sources to support them in their decision-making process. However, effective human-AI collaboration workflows are only beginning to emerge and remain particularly underexplored for inspection and SHM workflows. Similarly, the integration of immersive technologies with DT systems  have shown to further enhance user experience and decision-making. Combined, these technologies may offer novel pathways for inspection workflows remains largely unexamined. \par
%
The research addresses this gap in the following ways: 1) by building a system integrating DT, AI, and XR, 2) by developing and evaluating new methods for human--AI collaboration workflows with immersive DT for inspection, and 3) by providing a framework for future inspection workflows based on empirical evidence. \par
%
To increase the reliable use of AI for inspection workflows, it is necessary to evaluate how these are used by experts, how data needs to be presented, and how human--AI workflows can support the expert building situational awareness for decision-making. In the longer term, this project will provide the foundation on how DT, XR, and AI can be effectively combined and inform future workflows with AI that are reliable, safe, and easy to use. \par
%
%The technology around UAS, DTs, and AI are evolving rapidly. However, the application of these technologies remains largely incorporated into separate systems. The DT will act as the shared representation of the physical infrastructure, while a KG will organize the relationships among structural components, observations, inspection events, deterioration mechanisms, and engineering knowledge. AI will reason over this information and communicate its conclusions to engineers through traceable evidence rather than producing isolated recommendations. 
%
%\subsubsection{Scientific challenges}
%
%The project's distinguishing feature is not the use of UAS, DTs, KGs, or LLMs individually. These technologies are already developing rapidly. The novelty lies in integrating them within an explicit situational-awareness and human--AI decision framework in which AI reasoning and human engineering judgment have defined and complementary roles. The system will support evidence retrieval, contextualization, alternative interpretations, and identification of information gaps, while humans remain responsible for consequential structural and maintenance decisions.
%
%The main scientific focus will be on:
%
%\begin{enumerate}[label=F\arabic*]
%\item to cast an evolving semantic representation of the asset specific condition, historical record of deterioration and maintenance plans, and general knowledge on deterioration mechanisms.\label{focus:semantic}
%\item to develop an AI agent that utilizes such semantic representation to provide empirical evidence and formulate hypothesis on the construction condition and suggest maintenance plans.\label{focus:agentic}
%\item to create a human interaction layer based on XR and conversational agents through which the DT operators evaluate the evolving asset situation in an immersive environment with reduced cognitive load.\label{focus:immersive}
%\end{enumerate}
%
%AU-CAE focus is on \ref{focus:semantic} while ETHZ focus is on \ref{focus:immersive} whereas \ref{focus:agentic} is addressed jointly. 
%
The PhD project funded by \projectname will answer the following three questions aligned with the focus points:
%
\begin{enumerate}[label=Q\arabic*]
    \item what granularity is needed in the semantic representation of the asset to automate reasoning on its condition using AI agents? \label{question:sematic}
    \item how semantic representation shall be integrated with AI to maximize likelihood of logical consistency in reasoning?
    \item how can a conversational agent most effectively support the DT operator in providing relevant information? \label{question:interactive}
    \item how the DT operator shall interact with the immersive environment to improve situational awareness? \label{question:immersive}
\end{enumerate}
%
AU-CAE tackle \ref{question:sematic} while ETHZ will tackle \ref{question:interactive} and \ref{question:immersive}. The key-enabling technologies involved in this project are 1) semantic representations, 2) agentic AI 3) and XR providing the immersive layer. 

The PhD Candidate (Lorenzo Loyola) is educated as a structural engineer and hold +5 years of experience in R\&D (academia and industry) related to digitalization of structural design and monitoring workflow. He is already employed at AU. So the project will start in January 1st, 2027 and last 36m.
%
%how grounded AI can reason reliably across that representation, how to communicate uncertainty and alternative interpretations to engineers, and whether the resulting human--AI collaboration improves decision performance. 
%A further focus will be determining how much information is necessary for useful AI-supported SHM when complete DTs or extensive permanent sensing are unavailable.

\subsection{State of the Art}
%\subsubsection{Digital twins}

The implementation of SHM systems with physical sensing and methods for damage diagnosis and/or prognosis has required a great deal of work in system integration until recently. However, DT paradigm, emerged with the notable progress in cyber-physical system engineering, is changing this scenario\cite{talasila_digital_2026}. A DT combines digital and physical assets into a CPS to deliver specific services~\cite{fitzgerald_engineering_2024}. With the progress in IoT, cloud services, and edge computing, DTs for SHM today are capable of collecting multi-modal data (e.g., vibration and image data) in an almost automated manner~\cite{hoskere_unified_2025}. ML provides a workflow to automate specific processing tasks (e.g., crack localization). However, decision making regarding inspection and maintenance requires reasoning across large datasets and features, considering the entire history of the asset. This task is still a bottleneck that prevents us from exploiting the information contained in the data~\cite{celik_vision_2026}. Recent work points toward use of LLMs to fold DTs~\cite{luo_chattwin_2026}.
%A further limitation concerns transferability. Many advanced SHM and AI-supported infrastructure management approaches rely on relatively complete digital models, inspection histories, structured knowledge, imagery, or sensor networks. Their performance when information is incomplete, heterogeneous, or collected under constrained inspection resources remains insufficiently understood. 
%More recently, VLMs have emerged as a promising approach for image-based bridge damage documentation because of their ability to connect visual observations with natural-language descriptions and engineering information \cite{celik2026}.
%Sensors additionally provide continuous information on structural response and environmental conditions.
%This research addresses these gaps by integrating situational awareness, semantic representation, multimodal reasoning, human--AI interaction, and data sufficiency into a unified, experimentally validated SHM framework. The objective is therefore not only to improve the acquisition or automated interpretation of bridge information, but to investigate how heterogeneous and incomplete evidence can be transformed into traceable situational understanding that supports engineering decision-making.
Ontologies are currently being explored to add a semantic description to all components and functionalities of DTs to facilitate user interaction. An ontology is a formal description of the concepts in a particular area and how they relate to one another. In the context of SHM, ontologies are at the center of the current research agenda to represent relationships among structural components, inspection events, experimental observations, deterioration mechanisms, and maintenance actions \cite{ghidalia_combining_2024}.
%KGs are mathematical models utilized to describes \emph{ontologies} in computer programs. The integration of LLMs and KGs has similarly received growing attention as a means of enforcing explicit constraints on domain knowledge and semantic relationships.
%while natural-language interfaces such as \emph{ChatTwin} have demonstrated the potential of conversational interaction to improve accessibility to infrastructure DT information during operation and maintenance \cite{luo2026}.
Recent studies focusing on bridges demonstrated the potential of human-in-the-loop approaches, combining LLMs and ontologies to operate a DT using conversational agents in place of classical graphical user interfaces\cite{li_human---loop_2026}. 
%In parallel, LLMs and VLMs are beginning to support natural-language interaction with engineering information and multimodal evidence \cite{celik2026,luo2026}. 
%\subsubsection{Immersive environments}
Extended Reality (XR), which comprises immersive technologies such as Virtual Reality (VR), Augmented Reality (AR), and Mixed Reality (MR), shows great promise for enhancing DTs, both as a gateway to experiencing data that simulate or extend the real world~\cite{pandey2025digital} and as a means of improving usability~\cite{vona2025user}, collaboration and decision-making \cite{sthapit2026digital}. XR is also being explored as a way to enhance the manual bridge inspection process independently of DTs~\cite{smith_wearable_2022, gabbard_inspectar_2025}. Immersive DT environments have been combined with interactive virtual-agent interfaces for infrastructure leak simulation and monitoring \cite{marzban_immersive_2025}.
%Nevertheless, important gaps remain. 
The outlined state of the art is fragmented, and ongoing research is predominantly technology-centered, with comparatively limited attention to the engineer's situational awareness and decision-making process \cite{celik_vision_2026}. Human involvement is often introduced to validate AI outputs, refine generated knowledge, or facilitate interaction with digital information \cite{li_human---loop_2026,luo_chattwin_2026}, rather than being explicitly modeled as part of a joint human--AI reasoning and decision-making system. 

%Current approaches also provide limited evidence on how source provenance, uncertainty, competing hypotheses, and information gaps should be represented and communicated to engineers. Most importantly, relatively little empirical evidence exists on whether semantically grounded AI improves the quality, consistency, or efficiency of engineering decisions compared with conventional inspection workflows or generic AI assistance.
%These developments represent important steps toward more interactive DTs, but a broader challenge remains: moving beyond access to and visualization of information towards supporting engineers in interpreting how multiple observations relate to one another and how this evidence should inform subsequent inspection or maintenance actions.
%Graph-grounded AI approaches are therefore particularly promising because they offer a mechanism for constraining AI reasoning with domain-specific and asset-specific evidence rather than relying exclusively on the parametric knowledge of a language model.


%DT for monitoring of bridges (review) \cite{celik_vision_2026}%\url{https://scholar.google.com/citations?hl=en&user=LgEgwV4AAAAJ&view_op=list_works&sortby=pubdate}
%Combining ML and ontologies (review) \cite{ghidalia_combining_2024}
%Chat-with-twin is about using LLM to interact with DTs instead of GUIs \cite{luo_chattwin_2026}
%Virtual AI Assist. in immersive environment \cite{marzban_immersive_2025}
%AR to facilitate annotation and folding of damage state on structure \cite{gabbard_inspectar_nodate}
%LLM and GNN for ontology integration with AI \cite{li_human---loop_2026}

\subsection{Scientific Approach, Methodology, and Novelty}
\lengthnote{Describe the research contribution and potential impact }
% Add in the document:
\justifying
\setlength{\parindent}{0pt}
\setlength{\parskip}{0.8em}

The project will develop and experimentally validate a human-centered framework for bridge SHM in which AI and XR increase the operator awareness of the bridge condition. A \asset will be selected from the portfolio of structures available from FUBSE and UN. Reinforced concrete structures subjected to material deterioration (e.g., corrosion) will be prioritized. %work together to transform heterogeneous inspection data into traceable engineering knowledge and decision support. 
A central research challenge will be preventing the AI from presenting unsupported conclusions as established facts. A semantic description of the asset that capture the evolving condition of the bridge is used to constrain AI to provide engineering sound reasoning on the bridge state and collect evidence for decision support. The reasoning architecture will therefore distinguish between observed evidence, retrieved knowledge, inferred relationships, and hypotheses. %It will communicate uncertainty, provide evidence provenance, identify competing explanations, and explicitly indicate missing information. When available evidence is insufficient, the system will recommend what additional inspection, measurement, or historical information would most effectively reduce uncertainty.
%
The framework will represent bridge condition as an evolving infrastructure situation in which structural components, physical measurements, deterioration mechanisms, historical records, and maintenance activities are interconnected. As an example, the observation of a damage feature (e.g., a localized crack) will be recorded as contextualized empirical evidence describing \textit{what} was observed, \textit{when} it occurred, \textit{which} structural component it concerns. AI will be able to guess \textit{which sources} are likely to have produced the evidence, and its relationship with previous observations. This will allow the DT to evolve as new evidence becomes available and provide a machine-interpretable representation of the infrastructure situation.
%
The research program of \projectname is organized into three interconnected Work Packages (WPs) using \asset, which will be selected from a pool of alternatives in the global south supplied by the University of Salvador, Bahia, Brazil, and the University of Nairobi, Kenya. Particular attention will be given to water ingress and moisture-related deterioration because these phenomena are particularly relevant for aging infrastructures. Inspection data and maintenance decision scenarios will be formulated accordingly. 
%
WP1 develops the semantic representation of the physical asset and its evolving condition; WP2 develops multimodal, graph-grounded AI for evidence-based interpretation and reasoning; and WP3 develops the human--AI collaboration environment and evaluates how the resulting system affects engineering situational awareness and decision-making.

\textbf{WP1 Asset knowledge representation:} The main scientific contribution of WP1 will be a semantic model that links all digital assets associated with the \asset:\par
%
\textbf{T1.1 Digital assets:} Technical drawings, historical data, maintenance reports associated with the \asset will be gathered and dissected. Based on this input, a BIM and a FEM model of the \asset will be produced. The FEM will be used to describe the mechanical response of the bridge and to generate synthetic response data, whereas the BIM model will describe its geometry. Both available experimental response data and synthetic response data will be organized into databases (e.g., \texttt{InfluxDB} for time-series data and \texttt{MongoDB} for text data).\par
\textbf{T1.2 Semantic description:} A semantic model will be developed to describe the condition of the \asset, including the history of inspection data and maintenance reports. The BIM model developed in T1.1 will be used to derive a semantic layer for the \asset geometry. Additional semantic layers will be included for \emph{deterioration mechanisms} (e.g., cracking, corrosion), \emph{measurements} (e.g., UAS imagery, photogrammetric models), \emph{inspection reports}, \emph{maintenance plans}. The DT semantic description will be implemented as KGs.\par
%
\textbf{WP2 Multimodal and Graph-Grounded Artificial Intelligence:} The main scientific contribution of WP2 will be a multimodal, graph-grounded AI architecture for traceable infrastructure reasoning that moves beyond defect recognition toward evidence-based interpretation of bridge condition:\par
%
\textbf{T2.1 Automated annotation of deterioration:} WP2 will investigate how AI can transform heterogeneous observations into evidence-grounded interpretations of bridge condition. CV and VLM technologies will identify and contextualize candidate observations from UAS imagery, photographs, inspection documentation, and other multimodal information. The objective is not to develop another universal defect detector, but to determine how automatically extracted observations can become structured evidence within the semantic DT.\par
\textbf{T2.2 Automated reasoning on structure condition:} A graph-grounded LLM/VLM reasoning architecture will retrieve relevant asset-specific information from the KG before generating an engineering interpretation. For example, when moisture is detected below an expansion joint, the system may retrieve the joint's condition and inspection history, drainage configuration, recent rainfall exposure, nearby cracking or corrosion observations, bearing locations, previous repairs, and relevant engineering knowledge before proposing plausible deterioration hypotheses. The quality of AI-supported assessment changes as different information sources are removed or become incomplete will be also investigated. %This will establish whether reliable, trustworthy reasoning can be achieved in bridges where complete DTs, historical records, or extensive permanent sensor networks are unavailable.

\textbf{WP3 Development and Evaluation of an immersive Human--AI collaboration Workflows Framework:} WP3 will develop and evaluate methods of human-AI collaborative workflows that support trust, appropriate reliance, and decision-making for bridge inspection within immersive environments:\par
%
\textbf{T3.1 Human--AI Collaboration Workflows: }T3.1 develops and implements the conversational agent connected to the DT, building the basis of the human--AI workflow. It will Assoc. physical bridge components with current observations, historical changes, sensor information that can be queried by the user to build situational awareness.\textbf{ }\par
\textbf{T3.2 Immersive Component (XR): }T3.2. will enable immersive processes for experiencing the DT. Together with stakeholders, XR will be explored for their use with DTs for on-site and remote use, and paired with the workflow presented in T3.1. \par
\textbf{T3.3 Experimental Evaluation: }The complete framework will be evaluated experimentally using realistic bridge assessment scenarios in the lab, with user evaluations and in the field in partnership with stakeholders. Evaluation will examine decision accuracy and consistency, the time required to retrieve and interpret relevant evidence, situational awareness, cognitive workload, confidence calibration, appropriate reliance on AI, and the perceived usefulness and traceability of AI explanations. Human feedback will subsequently become part of the evolving semantic representation that will be refined iterately. \par

%WP3 will develop the human interaction layer through which engineers interrogate and evaluate the evolving infrastructure situation. Firstly, an interface using conversational agents connected to the DT will Assoc. physical bridge components with current observations, historical changes, sensor information, supporting evidence, AI-generated hypotheses, uncertainty, alternative explanations, and potential inspection or maintenance actions. 

%Second, the research will investigate XR as an interface for \textit{human--AI collaborative reasoning}. XR will be explored for their use with DTs in addition to the human interaction layer. Human feedback will subsequently become part of the evolving semantic representation, creating an iterative reasoning loop between the engineer, AI, and DT.

%Engineers will be able to inspect the evidence supporting an AI interpretation, interrogate relationships between observations, compare current and historical conditions, challenge AI-generated hypotheses, introduce expert knowledge, and reject unsupported conclusions. Human feedback will subsequently become part of the evolving semantic representation, creating an iterative reasoning loop between the engineer, AI, and DT.

%Engineering performance will be compared across three conditions: (1) conventional inspection and SHM workflows, (2) AI assistance without semantic grounding, and (3) the proposed graph-grounded human--AI framework. Evaluation will examine decision accuracy and consistency, time required to retrieve and interpret relevant evidence, situational awareness, cognitive workload, confidence calibration, appropriate reliance on AI, and the perceived usefulness and traceability of AI explanations.

%Field validation will include bridge case studies focusing particularly on water ingress, expansion-joint and drainage deficiencies, moisture exposure, cracking, and corrosion-related deterioration. A major international \asset in Salvador, Bahia, Brazil, will be conducted in collaboration with the Federal University of Bahia. The case will additionally support the data-sufficiency experiments by evaluating system performance under progressively different levels of information availability.

%WP3 will contribute to the scientific literature by assessing new methods of human-AI collaborative workflows that support trust, appropriate reliance, and decision-making for bridge inspection and SHM and present empirical evidence on their effectiveness. 

%whether semantic and graph-grounded human--AI collaboration improves engineering situational awareness and decision performance compared with conventional workflows and non-grounded AI assistance.

%\textbf{Integrated Scientific Contribution}
%
%Together, the three WPs establish a continuous research chain from infrastructure evidence to semantic representation, from semantic representation to grounded AI reasoning, and from AI reasoning to human engineering decisions. WP1 addresses how heterogeneous bridge information can be represented as an evolving infrastructure situation; WP2 investigates how AI can reason over that representation while maintaining provenance and uncertainty; and WP3 evaluates novel methods for human-AI collaboration workflows that aim to effectively calibrate trust and establish appropriate reliance that build the foundation for effective decision-making. 
%
The overall scientific contribution will be a new paradigm for an evidence-based decision-support system with immersive human-interaction facilitated by AI and AR. %supported by a semantic SHM representation, a multimodal graph-grounded reasoning architecture, an AR-enabled human-interaction framework, and experimental evidence on decision performance and data sufficiency. 
Practical impacts include faster interpretation of inspection information, reduced effort spent searching fragmented records, better prioritization of inspection resources, and more transparent and traceable documentation of scientific evidence in support to engineering decisions.

\subsection{Research Environment and Supervision}
\lengthnote{Description of research environment and infrastructure (demonstrating feasibility of the proposed project), Description of potential benefits and synergies expected through the interdisciplinary collaboration}

\justifying
\setlength{\parindent}{0pt}
\setlength{\parskip}{0.8em}
The PhD project is hosted at AU, Assoc. Prof. Giuseppe Abbiati is the main supervisor, and Assist. Prof. Jhonattan Martinez is the co-supervisor. The PhD Candidate (Lorenzo Loyola) will spend approximately 18 months at ETHZ under the supervision of Dr. Carolin Reichherzer and Prof. Ueli Angst.

%The project will be embedded in a highly interdisciplinary research environment jointly supported by AU and ETHZ, bringing together complementing expertise. AU will serve as the host institution, providing computational resources, DT and 3D modeling platforms, UAS and sensing technologies, and AI development tools. ETHZ expertise in material durability and guidance in human-centred engineering, experimental design, and XR, strengthening the project's scientific foundation.

%The project will be embedded in a highly interdisciplinary research environment jointly supported by Aarhus University (AU) and ETH Zurich, bringing together expertise in structural engineering, Structural Health Monitoring (SHM), infrastructure durability, DTs (DTs), Unmanned Aerial Systems (UAS)-based inspection, multimodal Artificial Intelligence (AI), and human-centered engineering. AU will serve as the host institution and provide the primary research environment, including computational resources, DT and 3D modeling platforms, UAS and photogrammetry capabilities, sensing technologies, graph-database infrastructure, and AI development tools required to develop and validate the proposed framework. ETH Zurich will complement this environment with internationally recognized expertise in infrastructure durability and corrosion science, strengthening the project's scientific foundation.

\textbf{Assoc. Prof. Giuseppe Abbiati (AU)} expertise is on DT engineering for SHM applications including both model-based and data-driven methods. He will supervised the implementation of the DT for the \asset.  %structural dynamics, computational methods, experimental testing, SHM, and infrastructure DTs. His role will provide the structural-engineering foundation for the project and ensure that the DT, monitoring methods, and AI-supported assessments remain connected to physically meaningful structural behavior. He will particularly contribute to defining SHM indicators, interpreting structural responses, developing realistic bridge-assessment scenarios, and validating the DT and monitoring components.

\textbf{Prof. Ueli Angst (ETHZ)} expertise in on material durability with focus on construction materials. He will supervise the formulation of maintenance plans, definition of inspection data and reasoning chain for deterioration diagnosis and prognosis for the \asset.
%will provide complementary expertise in infrastructure durability, reinforced-concrete deterioration, corrosion science, and service-life assessment. His contribution is particularly relevant to the project's focus on water ingress and moisture-related deterioration, where observed symptoms must be connected to physically plausible degradation mechanisms. His expertise will strengthen the definition of deterioration scenarios and the engineering knowledge represented within the semantic model, as well as the validation of AI-generated interpretations concerning moisture exposure, corrosion, and durability.

\textbf{Assist. Prof. Jhonattan G. Martinez (AU)} expertise is on automated inspections using UASs and computer-vision methods, BIM modeling. He will lead the formulation of the semantic representation for the DT and its integration with LLMs.
%will contribute expertise in DTs, UAS, Internet of Things (IoT) sensing, Building Information Modeling (BIM), data integration, and AI-supported decision systems for the built environment. His role will focus on integrating multimodal infrastructure information within the DT, developing UAS-based inspection and sensing workflows, and connecting semantic representations with AI-supported engineering decision-making. The complementary expertise of Abbiati and Martinez provides a strong internal AU supervisory environment connecting structural engineering and SHM with advanced digital technologies.

\textbf{Dr. Carolin Reichherzer (ETHZ)} expertise is in human-computer interaction in immersive environments and human-AI collaboration. She will lead the formulation of the human-centred immersive environment based on XR and its experimental validation.
%will provide expertise in the human-centred dimension, including human--AI collaboration, user experimental design, and XR. Her expertise will guide how engineers interact with AI-supported DTs, in both conventional and XR interfaces, particularly how interpretations, their supporting evidence, and their reliability should be presented to the engineer. She will also advise on the experimental design used to evaluate situational awareness, appropriate reliance, and decision quality, and on testing whether immersive interaction offers measurable benefits.

\textbf{Prof. Dayana Bastos (FUBSE)} expertise is on structural engineering and SHM. He will provide access to local infrastructures including possible field measurements in Salvador, Bahia, Brazil.

\textbf{Prof. Njeru Njue (UN) expertise is on geomatics and construction management. He will provide access to local infrastructures including possible field measurements in Nairobi, Kenya}
%\textbf{Federal University of Bahia (UFBA)} it will share access to the \asset and .  will contribute expertise in structural engineering and SHM and support the identification and development of bridge case studies . %This collaboration will provide access to a contrasting infrastructure context and enable the project to evaluate the proposed methodology under different climatic conditions, levels of data availability, inspection practices, and resource constraints.

\subsection{Stakeholders involved in the project}
\lengthnote{Description of which stakeholders will be involved and how.}
% Add in the preamble if not already included:


% ---------------------------------------------------------
% Project Partners and Stakeholders
% ---------------------------------------------------------

\justifying
\setlength{\parindent}{0pt}
\setlength{\parskip}{0.8em}

%The project will involve key academic and public-sector stakeholders throughout requirements definition, system development, field validation, and evaluation. 

%The project partners are:
%\begin{itemize}
%\item \textbf{Aarhus University, Department of Civil and Architectural Engineering (AU):} lead the project's scientific and technological development, including the DT (DT), Structural Health Monitoring (SHM) integration, semantic knowledge representation, multimodal Artificial Intelligence (AI) reasoning, Unmanned Aerial Systems (UAS)-based inspection, and human–AI evaluation. AU will also provide the primary research infrastructure and coordinate the interdisciplinary research activities.
%\item \textbf{ETH Zurich, Research Group of Durability of Engineering Materials}: will contribute complementary scientific expertise in infrastructure durability, reinforced-concrete deterioration, corrosion science, and service-life assessment. This expertise will strengthen the physical interpretation and validation of water- and moisture-related deterioration mechanisms represented within the semantic model and AI reasoning framework. ETH Zurich's contribution will therefore help ensure that AI-generated interpretations are grounded not only in semantic relationships and observational evidence, but also in physically meaningful deterioration mechanisms.
%\end{itemize}

Besides the project partners (AU and ETHZ) \projectname involves the following stakeholders:
%Additional stakeholders are:
\begin{itemize}
    \item \textbf{Federal University Of Bahia School Of Engineering (FUBSE), Department Of Structural And Construction Engineering, Brazil:} It will share access to local bridge portfolios for the selection of the \asset. Situated in a global south. %Through this, we will support the international validation of the framework through bridge and critical-infrastructure case studies in Salvador, Bahia. Prof. Dayana Bastos and collaborating researchers will contribute structural-engineering and SHM expertise, contextual knowledge, case-study development, and interpretation of local infrastructure conditions. The Brazilian cases will be particularly important for evaluating the methodology's applicability and transferability to developing-country contexts, including situations characterized by incomplete historical records, heterogeneous data availability, limited monitoring infrastructure, and constrained inspection and maintenance resources. global south.
    \item \textbf{University Of Nairobi (UN),  Department Of Real Estate, Construction Management And Quantity Surveying, Kenya:}  It will share access to local bridge portfolios for the selection of the \asset. Situated in a global south. 
    %\item \textbf{Aarhus Municipality:} It will share input on local infrastructure portfolios and will be exposed to the research progress. %will provide the essential connection to real infrastructure-management practice and support access to relevant infrastructure case studies in Aarhus. Its involvement will help define realistic inspection and maintenance scenarios, identify operational requirements and available data, and support validation of the proposed framework under actual infrastructure-management conditions. This collaboration will ensure that the research addresses practical challenges faced by public infrastructure owners rather than being limited to laboratory demonstrations. Danish public organization.
    %\item \textbf{COWI - SHM Heavy Civil:} Danish consulting company world leader in bridge engineering. It will be exposed to the project progress.
\end{itemize}

\pagemark{Project Description -- end}


% ==============================================================================
% 3. PROJECT RELEVANCE TO CEBE RESEARCH FIELDS  (max 1/2 page)
% ==============================================================================
\pagemark{CEBE relevance -- start}
\section{Project relevance to CEBE research fields}
\lengthnote{max 1/2 page (per the call)}
\lengthnote{Identification of the relevant research field(s) for the project Description of the expected contribution to the research field(s)}
\justifying
\setlength{\parindent}{0pt}
\setlength{\parskip}{0.8em}

%\projectname aligns with the scope of CEBE RF7 "Lifetime extension of the built environment" with a particular focus on bridge infrastructures and their impact on the global south. \par
%
%The ambition is to leverage the recent advancements in LLM-based AI to create DTs that enable the folding of large infrastructure portfolios, accompany inspections, and provide evidence for decision making with limited cognitive load and qualified personnel. A central component of \projectname is the human-AI interaction based on XR and AI and the methodological work necessary to convey expert know-how in such a DT. This specific challenge is not addressed by the RF7 baseline centered at AU.\par
%
%The project team specifically targeted the global south since underfunded inspection programs, along with climate conditions that favor material deterioration (high temperatures and humidity), undermine the resilience of urban communities.

%The project lies at the intersection of structural and civil engineering, digitalization of the built environment, AI, infrastructure resilience, and human-centered engineering systems. Its central contribution to CEBE is the development of an interdisciplinary methodology that transforms heterogeneous infrastructure observations into trustworthy and traceable engineering decision support. 
%1) Within structural engineering and SHM, the project advances the interpretation of deterioration by relating observed bridge conditions to structural components, inspection history, environmental exposure, and physically plausible deterioration mechanisms. Within DT research, it advances the role of DTs from primarily monitoring and visualization tools toward interactive, semantically grounded engineering reasoning environments. Within AI research for the built environment, the project investigates how multimodal LLM and VLM  systems can be grounded through KGs, infrastructure-specific evidence, uncertainty representation, and explicit human oversight.
%2) The project also contributes to human-centered digital engineering by investigating how reasoning responsibilities should be distributed between AI systems and structural engineers and by experimentally evaluating the performance of the combined human--AI system. This provides a stronger contribution than assessing algorithmic accuracy alone, because the relevant outcome in infrastructure management is not simply whether an AI model correctly identifies a defect, but whether the combined evidence, reasoning, and human judgment lead to more accurate, consistent, transparent, and timely engineering decisions.
%3) The collaboration with University of Salvador, Bahiha, Brazil and University of Nairobi, Kenya further strengthens the project's impact on global south across different climatic conditions. This international perspective will provide evidence regarding the framework's robustness, transferability, and applicability across both highly digitalized and resource-constrained infrastructure contexts.
%The expected contribution to CEBE is a transferable, human-centered model that integrates structural-engineering knowledge, semantic DTs, AI, and human expertise to support more resilient, intelligent, and evidence-based decision-making for the built environment.

\textbf{RF7 -- Extending the life of the built environment (primary).} \projectname addresses the RF7 mission of ``decision support that integrates structural health monitoring information''. Within WP7.1, it builds on the AU-led robot-assisted SHM systems by adding the layer that turns their observations into condition hypotheses linked to deterioration mechanisms such as corrosion (7.1.3). Within WP7.4, it contributes to risk-based decision-making for existing structures (7.4.3) by quantifying how engineers rely on AI-supported assessments. It thereby targets two barriers identified in the RF7 roadmap, limited confidence in monitoring data and new digital tools by practitioners.%, and provides evidence for the impact lever ``decision makers favour adaptation over demolition''. The key contribution is a human--AI decision process validated in terms of decision quality.%, not algorithmic accuracy alone.

\textbf{RF4 -- Digitalisation and automation.} The semantic DT and KG of WP1 extend the semantic modelling, reality capture and interoperability of WP4.4--4.6 from construction to the operation and maintenance phase. The project will contribute to the RF4/RF7 knowledge exchange on DT engineering and agentic AI identified in the RF7 roadmap.

\textbf{RF5 -- Climate resilient and adaptive infrastructure.} Water ingress is driven by precipitation and exposure. The KG will represent environmental exposure as context for interpreting moisture and corrosion symptoms, interfacing with RF5 work on precipitation data (5.1.4) and concrete corrosion mechanisms (5.4.1), and it will supply asset-level evidence to risk-based adaptation decisions (WP5.5).

\textbf{RF1 -- Assessing and measuring sustainability} The environmental consequences of case-study decisions will be screened with the RF1 sustainability framework (see Section~4).

\pagemark{CEBE relevance -- end}

% ==============================================================================
% 4. SUSTAINABILITY GOALS AND RELATION TO THE CEBE PROGRAM VISION  (max 1/2 page)
% ==============================================================================
\pagemark{Sustainability -- start}
\section{Sustainability goals of the project and relation to the CEBE program vision}
\lengthnote{max 1/2 page (per the call)}
\lengthnote{Specific sustainability objectives taking into account possible sustainability-related trade-offs, Sustainability-related conceptual and methodological choices}

\justifying
\setlength{\parindent}{0pt}
\setlength{\parskip}{0.8em}

%\textbf{Sustainable Cities and Communities (11):} Bridge infrastructures are critical to connect settlements and guarantee delivery of goods and services. Evidence-grounded decision support can help prioritize structural components requiring additional investigation or intervention while reducing unnecessary inspections and avoiding excessive permanent instrumentation while ensuring that critical infrastructures remain viable.

%\textbf{Climate Action (13):} Temperature raise and increase of rainfall negatively impact on corrosion mechanisms affecting common construction materials like concrete and steel. Underfunded inspection programs are likely to fail in tracking critical deterioration scenarions. \projectname aim to automate the process so that it can be deployed in large scale with limited qualified personnel. %Better identification of water ingress, drainage failures, moisture exposure, and associated deterioration may enable localized preventive interventions before damage develops into more material-intensive rehabilitation or premature replacement.

\projectname contributes to CEBE vision of a built environment in harmony with planetary boundaries by addressing risk-informed life extension (Gap 7, RF7) based on timely and trustworthy interpretation of deterioration. In bridges, much avoidable material use and CO$_2$ stems from late detection. Failing expansion joints and drainage let water and chlorides reach the reinforcement, and a local repair escalates into deck rehabilitation or replacement.

\textbf{Objectives.} (O1) \textit{Least intervention:} \projectname support decisions (inspect, repair, replace) that select the minimum action consistent with safety. T3.3 will record the intervention intensity of decisions alongside their correctness. (O2) \textit{Information sufficiency:} identify the minimum information needed for a trustworthy assessment, measurements and computation are added only where they change the decision (WP2). %(O3) \textit{Transferability:} show that the framework remains effective where expertise and inspection budgets are scarce, as in the Brazilian case study.

\textbf{Trade-offs and methodological choices.} Digitalization is not sustainable per se. LLM/VLM inference consumes energy and resources. \projectname therefore adopts a fit-for-purpose principle. Existing records are exploited before new data are acquired. Postponing action is sustainable only if deterioration is not underestimated. The separation of evidence, prior knowledge and inference, and engineer oversight are the safeguards against automation bias. \projectname supports, rather than presumes, life extension. %Its scope is the maintenance, repair and replacement stages (EN 15978, B2--B5), and the environmental consequences of case-study decisions will be screened with the RF1 sustainability framework.
%The project's primary sustainability objective is to support infrastructure life extension through earlier recognition, improved interpretation, and more targeted management of bridge deterioration.  The project therefore connects digital monitoring and evidence-based decision support with reduced material consumption, embodied carbon emissions, construction activity, and disruption over the infrastructure life cycle.

%A second objective is to improve the efficiency of inspection and maintenance resources. Infrastructure owners cannot inspect, monitor, or intervene across all components with equal intensity.  The proposed approach therefore seeks to direct limited financial, technological, and human resources toward the locations where they provide the greatest engineering value.

%The research explicitly recognizes the sustainability trade-offs associated with digitalization. DTs, UASs, sensors, AI, and data infrastructure themselves require energy, hardware, specialized skills, and financial resources. More digital technology is therefore not automatically more sustainable. The project will adopt a \textit{fit-for-purpose digitalization} principle and investigate the minimum combination of inspection, structural, environmental, and historical information required to achieve trustworthy AI-supported decision-making.

%The Brazilian case is particularly important for this methodological objective because it will evaluate whether the framework remains effective under incomplete and heterogeneous information conditions rather than only within highly instrumented infrastructure environments. In this way, the project addresses three interconnected dimensions of sustainability: \textbf{environmental sustainability} through infrastructure life extension and reduced demand for material-intensive interventions; \textbf{economic sustainability} through more efficient allocation of inspection and maintenance resources; and \textbf{social sustainability} through safer, more reliable, and transferable management of critical infrastructure across different technological and resource contexts.

\pagemark{Sustainability goals -- end}

% References for the Project Description (kept here since the call's
% 4-page limit for this section explicitly excludes the reference list).
\bibliographystyle{ieeetr}
\bibliography{application/references_ga, application/references}

\end{document}
